% ============================================================
\section*{Abstract}
\addcontentsline{toc}{section}{Abstract}
% ============================================================

Cosmic-Ray Neutron Sensing (CRNS) is a field technique for non-invasive,
area-averaged soil moisture measurement at the hectometre scale.
Before deploying a CRNS station, a thorough characterisation of the site
is necessary to (i)~quantify the geometric and geophysical corrections
that relate measured neutron counts to soil moisture, (ii)~estimate the
expected signal dynamic range, and (iii)~design the optimal soil sampling
campaign for calibration.

This technical note describes a fully integrated Python pipeline that
automates CRNS site characterisation from publicly available global
datasets.
The core physical outputs are three multiplicative correction factors:
the topographic correction $\kT$, computed by a three-dimensional
ray-casting algorithm over a 30-m DEM; the muon field-of-view correction
$\kM$, derived analytically from the horizon angle profile; and the
land-use/land-cover correction $\kL$, based on ESA WorldCover and
OpenStreetMap.
Their product $\kTot = \kT\,\kM\,\kL$ directly enters the calibration
of the reference count $\Nzero$ through the Desilets~(2010) equation,
extended here to account for clay-mineral lattice water $\lw$
(K\"ohli et al., 2021) while preserving full analytical invertibility.

Supplementary characterisation modules provide: seasonal atmospheric
water-vapour and snow-water-equivalent climatologies from ERA5;
above-ground biomass hydrogen estimates from MODIS LAI;
global lithological information from the Macrostrat API; an optimally
weighted soil sampling plan; and a Topographic Wetness Index.

The pipeline is validated on a test site in the Italian Alps, demonstrating
how each correction factor and supplementary quantity can inform
installation decisions and calibration planning.
